%% papericai.tex — Springer LNCS Proceedings Paper
%% Auto-generated from Master's Thesis
%% Target: 12–16 pages
%%
%% Thesis: "Accent Based Model for Speech Spoofing Countermeasures"
%% Author: Jaime Arturo Hurtado Romero
%% Advisor: Prof. Ruben Francisco Manrique
%% FLAG Lab, Universidad de los Andes, Bogotá, Colombia
%%
%% All figures are PDF; no external converter is required to build this file.

\documentclass{llncs}

\usepackage{graphicx}
\usepackage{amsmath}
\usepackage{amssymb}
\usepackage{multirow}
\usepackage{array}
\usepackage{bbding}
\usepackage{url}
 
\begin{document}

\title{Accent-Based Evaluation of Speech Anti-spoofing
       Countermeasures Across Multiple Languages}

%% The corresponding author is marked with \Envelope, as required by the
%% Springer proceedings instructions. This must match the name signing the
%% licence-to-publish agreement.
\author{Jaime Arturo Hurtado Romero\Envelope \and Tom\'{a}s Acosta Bernal \and
        Rub\'{e}n Francisco Manrique}

\institute{FLAG Lab, Department of Systems and Computing Engineering,\\
           Universidad de los Andes, Bogot\'{a}, Colombia\\
           \email{ja.hurtado@uniandes.edu.co}\\
           \email{\{t.acosta, rf.manrique\}@uniandes.edu.co}}

\maketitle

%% ----------------------------------------------------------------
\begin{abstract}
Speech-based biometric systems face growing threats from spoofing attacks
built on deep generative models, and countermeasures trained on English
corpora are known to degrade severely when applied to Spanish speech.  This
paper evaluates 48 anti-spoofing architectures, crossing four front-end
feature extractors (MFCC, LFCC, CQCC and spectrogram) with three Light
Convolutional Neural Network (LCNN) back-ends and four loss functions, over
the ASVspoof~2019 English database, the HABLA Spanish multi-accent database
and a combined bilingual corpus.  Each configuration is trained six times
under different random seeds and reported with medians, observed ranges and
distribution-free paired tests of Equal Error Rate (EER), which show that the
leading configurations are not statistically separable and should be read as
a family rather than a ranking.  Joint English-Spanish training reaches
median EER below $0.2\%$ on Spanish while sacrificing under one percentage
point on English, although a pool of conversion target identities shared
across partitions makes this an upper bound.  A per-attack decomposition
shows that the cross-language collapse is driven almost entirely by voice
conversion rather than by text-to-speech: English-trained models remain near
chance on converted speech yet still detect Spanish TTS attacks at $1.79\%$
EER.  A filter-bank scale study shows that descendant scaling benefits
Mel-scale front-ends but not linear ones, and an accent- and gender-resolved
analysis over five Latin-American varieties identifies speaker gender rather
than regional accent as the dominant residual source of error.

\keywords{speech anti-spoofing \and accent diversity \and cross-language
          evaluation \and voice conversion \and LCNN \and ASVspoof \and HABLA}
\end{abstract}

%% ----------------------------------------------------------------
\section{Introduction}

Biometric systems that rely on speech for identity verification are exposed
to a growing range of spoofing attacks.  \textit{Speech synthesis}
(text-to-speech, TTS) produces spoken output from written input, while
\textit{voice conversion} (VC) modifies non-linguistic characteristics of a
waveform to impersonate a target speaker~\cite{stylianou1998continuous}.
Both attack types can be mounted at the logical access point of an automatic
speaker verification (ASV) system and may lead to unauthorised access,
financial fraud, or identity theft~\cite{xue2023physiological}.  The
proliferation of Generative Adversarial Networks (GANs)~\cite{goodfellow2020generative}
and Variational Autoencoders (VAEs)~\cite{kingma2013auto} has dramatically
lowered the barrier to generating perceptually convincing synthetic speech,
underscoring the urgency of robust anti-spoofing countermeasures
(CMs)~\cite{kaur2014review}.

The \textit{ASVspoof} challenge series~\cite{yamagishi2019asvspoof,delgado2021asvspoof} has driven most recent progress in anti-spoofing;
however, all its benchmark datasets are exclusively in English.
Prior work has shown that state-of-the-art CMs trained on English data
degrade severely when evaluated on Spanish speech~\cite{tamayo2022voice},
with Equal Error Rates exceeding $30\%$ even for strong baselines.  This
degradation is linked to accent-driven differences in vowel space, formant
structure, and speaking rate that directly affect the cepstral features
commonly employed as front-ends.  Despite this, there is limited research
that \textit{systematically} quantifies how specific front-end/back-end
combinations respond to language and accent variation.

This work makes the following contributions:

\begin{enumerate}
  \item A systematic comparison of four front-end feature extractors combined
        with three LCNN back-end variants and four loss functions (48
        architecture combinations), each trained and evaluated across three
        dataset conditions and reported with distribution-free statistics
        over six random seeds, so that architectural differences can be
        separated from seed variation.
  \item A per-attack decomposition of the cross-language degradation
        reported by Tamayo et~al.~\cite{tamayo2022voice}, showing that the
        collapse is driven almost entirely by voice conversion rather than by
        text-to-speech, a distinction not previously reported.
  \item An accent and gender-resolved analysis over five Spanish regional
        varieties (Argentina, Chile, Colombia, Peru, Venezuela), identifying
        speaker gender as a residual error source whose direction reverses
        between countries.
  \item A filter-bank scale analysis showing that descendant
        (low-frequency emphasis) scaling benefits Mel-scale cepstral
        front-ends but not linear-scale ones, refining the assumption that
        low-frequency emphasis is universally preferable.
\end{enumerate}

%% ----------------------------------------------------------------
\section{Related Work}

\subsection{Anti-spoofing Datasets}

The ASVspoof benchmarks are the primary evaluation corpora for logical-access
(LA) and physical-access (PA) spoofing.  ASVspoof~2019~\cite{yamagishi2019asvspoof}
covers LA and PA attacks, employing 19 TTS/VC systems in training and 48 in
evaluation; ASVspoof~2021~\cite{delgado2021asvspoof} further extends to
deepfake (DF) speech.  All ASVspoof datasets are English only.  The
\textit{HABLA Spoof} corpus~\cite{tamayo2022voice} addresses the Spanish
gap: it contains synthesised and converted speech from speakers representing
five Latin-American accents, enabling accent-granularity evaluation of
anti-spoofing systems.  Table~\ref{table:datasets-taxonomy} summarises the
main datasets considered in this study.

\begin{table}[htbp]
\centering
\caption{Anti-spoofing dataset taxonomy. LA = Logical Access, PA = Physical
         Access, DF = Deepfake; EN = English, ES = Spanish.}
\label{table:datasets-taxonomy}
\setlength\extrarowheight{-2pt}
\begin{tabular}{lcl}
\hline
\textbf{Dataset}                           & \textbf{Lang.} & \textbf{Attacks} \\
\hline
HABLA Spoof~\cite{tamayo2022voice}         & ES & LA, DF       \\
ASVspoof 2019~\cite{yamagishi2019asvspoof} & EN & LA, PA       \\
ASVspoof 2021~\cite{delgado2021asvspoof}   & EN & LA, PA, DF   \\
SAS~\cite{wu2016anti}                      & EN & LA           \\
FoR~\cite{reimao2019dataset}               & EN & LA           \\
VCTK~\cite{yamagishi2019cstr}              & EN & LA           \\
\hline
\end{tabular}
\end{table}

\subsection{Front-End Feature Extractors}

Digital Signal Processing (DSP)-based front-ends extract a compact acoustic
feature sequence from each utterance.  MFCC~\cite{davis1980comparison}
applies a Mel-scale filter bank followed by a Discrete Cosine Transform
(DCT), producing perceptually motivated coefficients compact in the low
frequencies.  LFCC~\cite{zhou2011linear} replaces the Mel scale with a
uniform linear-scale filter bank, treating all frequencies equally.
CQCC applies the Constant-Q Transform before cepstral analysis, providing
geometrically spaced frequency resolution suited to capturing spectral
artefacts in synthesised speech.  Spectrogram features retain the full
Short-Time Fourier Transform (STFT) magnitude, serving as the richest
but highest-dimensional representation.  Both LFCC and CQCC were the
official ASVspoof baselines~\cite{wang2022practical}.

\subsection{Back-End Classifiers}

The back-end maps the front-end feature sequence to a spoofing score.
Light Convolutional Neural Networks (LCNNs)~\cite{lavrentyeva2017audio}
use Max-Feature-Map (MFM) activations to suppress irrelevant feature
channels and have become a strong standard baseline.
LSTM-augmented LCNNs~\cite{wang2021comparative} further model long-range
temporal dependencies.  End-to-end models such as RawNet2~\cite{tak2021end}
operate directly on raw waveforms, whereas classical Gaussian Mixture Models
(GMMs) remain competitive for PA scenarios~\cite{cai2017countermeasures}.
Representative front-end/back-end pairings from the literature span
spectrogram inputs with CNN, RNN and squeeze-excitation back-ends, cepstral
inputs with GMM and TDNN, and end-to-end models operating on raw waveforms.

\subsection{Research Gap}

Existing anti-spoofing literature is concentrated on English corpora and
single-language settings.  Tamayo~et~al.~\cite{tamayo2022voice,tamayoflorez23_interspeech} demonstrated
the cross-lingual drop but did not systematically evaluate front-end/back-end
combinations or filter-bank scales.  The effect of regional accents on
feature extraction and the degree to which multilingual training mitigates
language mismatch remain largely open questions, both of which our work
directly addresses.

%% ----------------------------------------------------------------
\section{Methodology}

\subsection{Experimental Framework}

Each anti-spoofing model is composed of a \textit{front-end} that extracts
acoustic features and a \textit{back-end} that classifies the utterance as
bona fide or spoofed.  We systematically combine all front-end and back-end
variants in Table~\ref{table:antispoofing-modules}, yielding $4 \times 3 \times
4 = 48$ configurations per training condition.

\begin{table}[htbp]
\centering
\caption{Anti-spoofing architecture modules explored in this study.}
\label{table:antispoofing-modules}
\setlength\extrarowheight{-2pt}
\begin{tabular}{lll}
\hline
\textbf{Front-End} & \textbf{Back-End}    & \textbf{Loss / Activation} \\
\hline
CQCC               & LCNN-Attention       & MSE for P2S-Grad \\
MFCC               & LCNN-trim-pad        & AM-Softmax       \\
LFCC               & LCNN-LSTM-sum        & OC-Softmax       \\
Spectrogram        &  ---                 & Sigmoid          \\
\hline
\end{tabular}
\end{table}

\subsection{Front-End Feature Extraction}

\textbf{MFCC and LFCC} are cepstral coefficient representations computed
frame-by-frame via Discrete Fourier Transform (DFT), log-magnitude, and DCT.
MFCC uses a Mel-scale (logarithmically spaced) filter bank that mimics human
auditory sensitivity; LFCC uses a linearly spaced bank that weights all
frequencies equally~\cite{zhou2011linear}.  We vary the filter-bank
\textit{scale response} in three modes that modulate how gain changes across
frequency: \textit{ascendant} (gain increases toward high frequencies),
\textit{constant} (flat gain), and \textit{descendant} (gain decreases toward
high frequencies, emphasising low-frequency formants).

\textbf{CQCC} applies the Constant-Q Transform before cepstral analysis,
providing finer frequency resolution at low frequencies where anti-spoofing
artefacts are often present.

\textbf{Spectrogram} retains the full STFT magnitude as a 2-D image fed
directly into the LCNN.

All front-ends share the same short-time analysis, so that differences in
performance are attributable to the filter bank and cepstral stage rather than
to framing: a Hann window of 320 samples with a 160-sample shift and a
512-point FFT.  MFCC, LFCC and CQCC retain 20 static coefficients and append
first- and second-order deltas; the spectrogram front-end uses 60 static
coefficients and no deltas.  All four therefore present 60 dimensions per
frame to the back-end.

No voice activity detection, cepstral mean normalisation or feature-level
normalisation is applied at any stage.

\subsection{LCNN Back-End Architecture}

All back-ends share a common LCNN feature-extraction stem followed by one of
three utterance-level pooling strategies.  \textit{Attention-based} pooling
learns to weight frames by relevance; \textit{trim-pad} forces fixed-length
input; \textit{LSTM-sum} models temporal dynamics with a bidirectional LSTM
before summing.  The pooled embedding is fed to a classification head whose
output distribution is shaped by the chosen loss function: AM-Softmax,
OC-Softmax, MSE-based P2S-Grad, or Sigmoid cross-entropy.

The LCNN stem architecture, detailed in Table~\ref{table:lcnn-architecture},
employs MFM activations after every convolution, batch normalisation where
indicated, and max-pooling for spatial downsampling.

\begin{table}[htbp]
\centering
\caption{LCNN stem architecture shared across all back-end variants.
         MFM = Max-Feature-Map, BN = Batch Normalisation.}
\label{table:lcnn-architecture}
\setlength\extrarowheight{-4pt}
\footnotesize
\begin{tabular}{lcc}
\hline
\textbf{Layer}         & \textbf{Filter/Stride}      & \textbf{Output}          \\
\hline
Conv + MFM             & $5{\times}5$ / $1{\times}1$ & $864{\times}600{\times}32$ \\
MaxPool                & $2{\times}2$ / $2{\times}2$ & $432{\times}300{\times}32$ \\
\hline
Conv + MFM + BN        & $1{\times}1$ / $1{\times}1$ & $432{\times}300{\times}32$ \\
Conv + MFM             & $3{\times}3$ / $1{\times}1$ & $432{\times}300{\times}48$ \\
MaxPool + BN           & $2{\times}2$ / $2{\times}2$ & $216{\times}150{\times}48$ \\
\hline
Conv + MFM + BN        & $1{\times}1$ / $1{\times}1$ & $216{\times}150{\times}48$ \\
Conv + MFM             & $3{\times}3$ / $1{\times}1$ & $216{\times}150{\times}64$ \\
MaxPool                & $2{\times}2$ / $2{\times}2$ & $108{\times}75{\times}64$  \\
\hline
4$\times$(Conv+MFM+BN) & $1{\times}1$--$3{\times}3$ / $1{\times}1$ & $108{\times}75{\times}32$ \\
MaxPool                & $2{\times}2$ / $2{\times}2$ & $54{\times}38{\times}32$   \\
\hline
\end{tabular}
\end{table}

\subsection{Datasets and Partitions}
\label{sec:datasets}

Two anti-spoofing corpora are used; partition sizes are listed in
Table~\ref{table:datasets-partitions}.

\textbf{English (ASVspoof~2019 LA)}~\cite{yamagishi2019asvspoof}: Standard
logical-access split with 19 TTS/VC systems in training data and 48 in
the evaluation set.

\textbf{Spanish (HABLA Spoof)}~\cite{tamayo2022voice,tamayoflorez23_interspeech}: Latin-American
multi-accent corpus covering five countries (Argentina, Chile, Colombia,
Peru, Venezuela).  We audited the partition assignment directly from the file
lists.  Bonafide speakers are strictly disjoint across the three splits (29,
16 and 14, none shared), as are the source speakers feeding the conversion
systems; the \emph{target} identities those systems imitate are not, with all
113 test identities also present in training.  No test speaker's genuine voice
is therefore seen during training, but conversions directed at the same
identities are, and this affects $95.5\%$ of Spanish test utterances.  Working
in the opposite direction, CycleGAN conversions appear in training but not in
the Spanish test split, making part of the evaluation an unseen-attack
condition.

\textbf{Combined}: The English and Spanish training sets are concatenated
to form a 60\,064-utterance bilingual training corpus; evaluation is
performed separately on each language's test partition.

Three properties of this design constrain the bilingual results.  First, the
merged corpus is \emph{not} language-balanced: Spanish contributes $34\,684$
utterances ($57.7\%$) against $25\,380$ for English ($42.3\%$), so the
near-zero Spanish EER and the small English penalty are confounded with corpus
composition and we do not claim to have separated the two.  Second, the
Spanish test partition holds only $4\,809$ utterances, so an EER of $0.1\%$
reflects a handful of trials, which motivates the treatment in
Section~\ref{sec:metric}.  Third, the shared pool of conversion targets leaves
the spoofed side of the decision exposed even though the bonafide side is not:
a model may learn artefacts specific to a target identity and meet that
identity again at test time.  The Spanish rates are therefore an upper bound
on performance under a strictly identity-disjoint protocol, while comparisons
across architectures, which share this exposure, remain internally valid.

\begin{table}[htbp]
\centering
\caption{Dataset partition sizes (number of utterances).}
\label{table:datasets-partitions}
\setlength\extrarowheight{-2pt}
\begin{tabular}{lccc}
\hline
\textbf{Training set} & \textbf{Train} & \textbf{Dev} & \textbf{Test} \\
\hline
English               & 25\,380 & 24\,844 & 71\,237 \\
Spanish               & 34\,684 &  3\,658 &  4\,809 \\
Combined              & 60\,064 & 28\,502 & --- \\
\hline
\end{tabular}
\end{table}

\subsection{Training Protocol}

All models are trained for up to 100 epochs with an Adam
optimiser~\cite{wang2021comparative} ($\beta_1=0.9$, $\beta_2=0.999$,
$\epsilon=10^{-8}$).  The initial learning rate of $3\times10^{-4}$ is
halved every ten epochs.  Mini-batches of 64 utterances are grouped by
similar duration.  Early stopping terminates training after 50 epochs
without improvement on the development loss; the checkpoint with the
lowest development loss is retained.  No voice activity detection or
feature normalisation is applied.  To account for variance due to random
initialisation, every configuration is trained \textbf{six} times with
different random seeds; results are reported as \emph{mean $\pm$ standard
deviation EER~(\%)}.

In a representative run, early stopping triggers at epoch~72, fifty epochs
after the last improvement in development loss.

\subsection{Evaluation Metric}
\label{sec:metric}

Models are evaluated using the \textit{Equal Error Rate} (EER), the
operating point at which the False Rejection Rate (FRR) equals the False
Acceptance Rate (FAR):

\begin{equation}
    \text{EER} = \frac{\text{FRR}(\tau') + \text{FAR}(\tau')}{2},
    \qquad
    \tau' = \underset{\tau}{\arg\min}\,|\text{FRR}(\tau) - \text{FAR}(\tau)|.
    \label{eq:eer}
\end{equation}

A lower EER indicates better discrimination between bona-fide and
spoofed utterances.  Because EER alone ignores the cost of the errors it
counts, we additionally report the minimum normalised tandem detection cost
function (min t-DCF), which weights countermeasure decisions by their effect
on a downstream speaker verification system, for both languages.

\subsubsection{Statistical treatment.}
The same six seeds are used for every configuration, so runs can be compared
in matched pairs.  EER is bounded below by zero and, in the Spanish
conditions, concentrates near that bound: exact zeros occur, skewness exceeds
one in a quarter of the cells, and a Gaussian interval would place mass on
negative error rates.  We therefore avoid any normality assumption, summarise
with the median and observed range, compare with the Wilcoxon signed-rank
test on seed-matched pairs, and report Cohen's $d$ as effect size, applying a
Bonferroni correction for multiple comparisons.  With six seeds the smallest
attainable two-sided $p$-value is $0.031$, so the design is underpowered for
small effects and we do not treat marginal $p$-values as decisive.

%% ----------------------------------------------------------------
\section{Experiments and Results}

We structure the experiments around five conditions, systematically varying
the training language and the evaluation language.  All models are
implemented on the open-source anti-spoofing framework released by the
Yamagishi Laboratory~\cite{Yamagishilab}, and the scripts used to audit the
corpus partitions reported in Section~\ref{sec:datasets} are included with
this submission.

\subsection{English Training, English Evaluation}

Table~\ref{table:mean-train-en-eval-en} reports results when all models
are trained and evaluated on the English partitions.  Performance is
consistently low across all 48 configurations.  Cepstral front-ends (LFCC and
MFCC) consistently outperform CQCC and Spectrogram.

The lowest mean EER is obtained by \textbf{LFCC-LCNN-LSTM-sum with P2S-Grad}
at $2.67\%$ ($\sigma=0.60$), but this configuration is not separable from its
competitors: against the next seven, seed-matched Wilcoxon tests yield no
significant difference in any case ($p > 0.31$ throughout) and the effect size
against the runner-up is $d = 0.22$.  Summarised without a normality
assumption, it has a median EER of $2.51\%$ over an observed range of
$2.12$--$3.70$, while the seven that follow span medians of $2.61$ to
$2.97\%$ with ranges that overlap it in every case, and their median min
t-DCF values ($0.063$ to $0.075$) do not reproduce the EER ordering.  The
eight should therefore be read as one indistinguishable group rather than a
ranking.  What the data does support at this sample size is that P2S-Grad and
OC-Softmax outperform AM-Softmax and Sigmoid, that Spectrogram front-ends show
higher seed variance than cepstral ones, and that CQCC performs comparably to
Spectrogram here.

\begin{table}[htbp]
\centering
\caption{Mean and std of EER (\%) for English train / English eval.
         Best per back-end in \textbf{bold}; overall best
         \underline{\textbf{underlined}}.}
\label{table:mean-train-en-eval-en}
\setlength\extrarowheight{-8pt}
\footnotesize
\begin{tabular}{llcccccc}
\hline
 & & \multicolumn{2}{c}{Attention} & \multicolumn{2}{c}{Trim-pad}
   & \multicolumn{2}{c}{LSTM-sum} \\
\textbf{Front} & \textbf{Activation} &
  \textbf{Mn} & \textbf{Sd} & \textbf{Mn} & \textbf{Sd} &
  \textbf{Mn} & \textbf{Sd} \\
\hline
\textbf{CQCC}
  & AM-softmax & 7.58 & 0.69 & 7.49 & 0.73 & 7.81 & 0.73 \\
  & P2SGrad    & 6.65 & 0.33 & \textbf{6.00} & \textbf{0.19} & 6.20 & 0.14 \\
  & OC-softmax & \textbf{6.40} & \textbf{0.30} & 6.12 & 0.12 & 6.63 & 0.32 \\
  & Sigmoid    & 8.08 & 1.01 & 6.23 & 0.32 & \textbf{6.14} & \textbf{0.22} \\
\hline
\textbf{Spec.}
  & AM-softmax & 5.35 & 1.53 & 5.52 & 0.56 & 4.67 & 0.76 \\
  & P2SGrad    & 5.68 & 0.71 & \textbf{3.79} & \textbf{0.71} & \textbf{3.48} & \textbf{0.82} \\
  & OC-softmax & 5.67 & 1.59 & 5.10 & 0.70 & 4.13 & 0.83 \\
  & Sigmoid    & \textbf{4.81} & \textbf{0.98} & 4.50 & 0.96 & 4.19 & 0.79 \\
\hline
\textbf{MFCC}
  & AM-softmax & 4.46 & 0.37 & 6.97 & 1.95 & 4.73 & 2.30 \\
  & P2SGrad    & 3.04 & 0.51 & \textbf{2.88} & \textbf{0.34} & \textbf{2.83} & \textbf{0.29} \\
  & OC-softmax & 3.71 & 1.37 & 4.49 & 1.09 & 3.51 & 0.93 \\
  & Sigmoid    & \textbf{3.50} & \textbf{0.17} & 3.52 & 0.61 & 3.47 & 0.94 \\
\hline
\textbf{LFCC}
  & AM-softmax & 4.00 & 0.33 & 5.58 & 0.61 & 3.25 & 1.11 \\
  & P2SGrad    & 3.48 & 0.78 & 2.88 & 0.41 & \underline{\textbf{2.67}} & \underline{\textbf{0.60}} \\
  & OC-softmax & \textbf{2.96} & \textbf{0.35} & \textbf{2.88} & \textbf{0.23} & 2.80 & 0.61 \\
  & Sigmoid    & 3.73 & 0.35 & 3.04 & 0.36 & 3.33 & 0.31 \\
\hline
\end{tabular}
\end{table}

\subsection{Cross-Language Evaluation: English Models on Spanish}
\label{sec:cross-language}

Applying the English-trained models directly to the Spanish test partition
reveals severe degradation: \emph{all} 48 configurations exceed $30\%$ EER,
the lowest configuration mean being $34.4\%$, and 28 of the 48 exceed $40\%$.
Among front-ends, Spectrogram degrades the least (best mean $\approx 36\%$).
CQCC and LFCC degrade the most, at $45.5\%$ and $45.4\%$ mean EER
respectively; the two are effectively tied and we do not read an ordering
between them.  MFCC and LFCC, despite being strong English baselines, yield
cross-lingual EERs of $34\text{--}55\%$.  This confirms the degradation
reported by Tamayo et~al.~\cite{tamayo2022voice}.

The aggregate figure conceals a structure that matters for interpretation.
Decomposing the same evaluation by attack family
(Table~\ref{table:per-attack}) shows that the collapse is not uniform: the
three voice conversion families drive it almost entirely, while
text-to-speech attacks remain detectable.  The best English-trained model
reaches $1.79\%$ EER on Spanish TTS attacks, within the range it achieves on
English, yet is at or below chance on CycleGAN, DiffVC and StarGAN.

\begin{table}[htbp]
\centering
\caption{EER (\%) by attack family on the Spanish evaluation set, for the
         three P2S-Grad back-ends.  English-trained models fail on voice
         conversion but not on TTS.  VC = voice conversion.}
\label{table:per-attack}
\setlength\extrarowheight{-2pt}
\footnotesize
\begin{tabular}{llcccc}
\hline
 & & \multicolumn{3}{c}{\textbf{VC attacks}} & \textbf{TTS} \\
\textbf{Back-end} & \textbf{Training} & CycleGAN & DiffVC & StarGAN & \\
\hline
LCNN-Attention & English & 53.09 & 50.33 & 64.77 & 46.11 \\
LCNN-LSTM-sum  & English & 51.72 & 53.99 & 65.83 & \textbf{1.79} \\
LCNN-trim-pad  & English & 51.28 & 43.22 & 63.77 & 16.92 \\
\hline
LCNN-Attention & Spanish & 0.03 & 0.00 & 0.01 & 0.00 \\
LCNN-LSTM-sum  & Spanish & 0.01 & 0.01 & 0.01 & 0.00 \\
LCNN-trim-pad  & Spanish & 0.03 & 0.01 & 0.01 & 0.01 \\
\hline
\end{tabular}
\end{table}

Two details are worth noting.  StarGAN yields EER \emph{above} chance
($63\text{--}66\%$), so the decision is inverted rather than merely
uninformative, and the TTS column varies by an order of magnitude across
back-ends, so cross-language TTS transfer depends on the pooling strategy
and not on the front-end alone.

A plausible reading is that vocoder artefacts are largely
language-independent and therefore transfer, whereas conversion artefacts are
entangled with the speaker and accent characteristics that the conversion
model preserves, and those are exactly what changes between English and
Spanish.  Figure~\ref{fig:det} shows the corresponding detection error
trade-off curves.

\begin{figure}[htbp]
    \centering
    \includegraphics[width=0.62\textwidth]{figures/det_curves}
    \caption{Detection error trade-off curves on the Spanish evaluation set.
             Grey curves are English-trained models, black curves are
             Spanish-trained models of the same architecture; circles mark the
             equal error rate.  Curves are computed on the full HABLA
             evaluation set (32{,}327 utterances), which is larger than the
             partition used for Table~\ref{table:datasets-partitions}.}
    \label{fig:det}
\end{figure}

\subsection{Spanish Training, Spanish Evaluation}

Training and evaluating exclusively on Spanish data yields low EER
for most architectures (Table~\ref{table:mean-train-es-eval-es}).  The best
configuration is \textbf{LFCC-LCNN-Attention with P2S-Grad} at $0.14\%$ EER
($\sigma=0.09$).  Spectrogram front-ends are the notable exception,
consistently yielding EER of $12\text{--}20\%$, suggesting that synthetic
Spanish speech lacks the texture-level discriminability that Spectrogram
features exploit in English.

Comparing Table~\ref{table:mean-train-es-eval-es} with the English baseline,
the \emph{architecture ranking shifts}: LFCC-LCNN-Attention/P2S-Grad is
best for Spanish whereas LFCC-LCNN-LSTM-sum/P2S-Grad was best for English.
This language-specific sensitivity motivates the combined-training approach.

\begin{table}[htbp]
\centering
\caption{Mean and std of EER (\%) for Spanish train / Spanish eval (best
         activation per front-end/back-end shown).  Overall best
         \underline{\textbf{underlined}}.}
\label{table:mean-train-es-eval-es}
\setlength\extrarowheight{-8pt}
\footnotesize
\begin{tabular}{llcccccc}
\hline
 & & \multicolumn{2}{c}{Attention} & \multicolumn{2}{c}{Trim-pad}
   & \multicolumn{2}{c}{LSTM-sum} \\
\textbf{Front} & \textbf{Activation} &
  \textbf{Mn} & \textbf{Sd} & \textbf{Mn} & \textbf{Sd} &
  \textbf{Mn} & \textbf{Sd} \\
\hline
\textbf{CQCC}
  & AM-softmax & 1.35 & 1.16 & 4.15 & 2.15 & 2.54 & 1.34 \\
  & P2SGrad    & \textbf{0.61} & \textbf{0.18} & 3.87 & 1.82 & 2.69 & 1.27 \\
  & OC-softmax & 2.69 & 2.15 & 4.73 & 1.87 & \textbf{1.29} & \textbf{0.48} \\
  & Sigmoid    & 1.29 & 0.63 & \textbf{3.36} & \textbf{1.93} & 2.53 & 1.43 \\
\hline
\textbf{Spec.}
  & AM-softmax & 17.15 & 3.33 & 17.98 & 4.05 & 17.30 & 1.50 \\
  & P2SGrad    & \textbf{12.14} & \textbf{3.02} & \textbf{17.24} & \textbf{1.43} & 19.50 & 4.96 \\
  & OC-softmax & 15.83 & 2.50 & 17.41 & 1.83 & \textbf{15.76} & \textbf{1.90} \\
  & Sigmoid    & 20.33 & 3.61 & 20.23 & 3.50 & 19.58 & 2.55 \\
\hline
\textbf{MFCC}
  & AM-softmax & 1.10 & 1.28 & 3.79 & 3.11 & 0.65 & 0.57 \\
  & P2SGrad    & \textbf{0.44} & \textbf{0.77} & \textbf{3.10} & \textbf{1.84} & 1.21 & 1.20 \\
  & OC-softmax & 0.77 & 0.73 & 3.42 & 4.82 & 1.38 & 1.14 \\
  & Sigmoid    & 1.74 & 2.34 & 3.25 & 2.83 & \textbf{0.42} & \textbf{0.31} \\
\hline
\textbf{LFCC}
  & AM-softmax & 1.40 & 0.90 & 3.58 & 3.05 & \textbf{0.45} & \textbf{0.27} \\
  & P2SGrad    & \underline{\textbf{0.14}} & \underline{\textbf{0.09}} & 3.55 & 2.80 & 2.32 & 1.57 \\
  & OC-softmax & 0.68 & 0.90 & 3.89 & 4.70 & 1.30 & 1.71 \\
  & Sigmoid    & 0.49 & 0.50 & 4.07 & 5.36 & 0.57 & 0.55 \\
\hline
\end{tabular}
\end{table}

\subsection{Combined Bilingual Training}
\label{sec:combined}

Training on the merged English-Spanish corpus provides the most practically
robust models.  Table~\ref{table:mean-train-un-eval-es} shows results on the
Spanish evaluation partition; Table~\ref{table:mean-train-un-eval-en} shows results
on the English partition.

MFCC and LFCC front-ends achieve near-zero EER ($\leq 0.5\%$) on Spanish
evaluation.  The lowest mean is \textbf{MFCC-LCNN-LSTM-sum with AM-Softmax}
at $0.087\%$ ($\sigma=0.04$), but here too the leading group is not
separable: against the two nearest competitors (both $0.17\%$) Wilcoxon
returns $p = 0.0625$, and although a paired $t$-test flags three of seven at
$p < 0.05$, none survives Bonferroni correction ($p < 0.0071$ required).  We
therefore report a leading family of four cepstral configurations below
$0.2\%$ rather than a single best system, whose median min t-DCF lies between
$0.0022$ and $0.0035$.  CQCC also performs well ($\leq 1.3\%$); Spectrogram
remains the weakest front-end, mirroring the Spanish-only experiment.

On the English evaluation partition (Table~\ref{table:mean-train-un-eval-en}),
performance decreases by only $\sim$1 percentage point relative to the
English-only baseline, the lowest bilingual English mean being
MFCC-LCNN-trim-pad with P2S-Grad at $3.64\%$.  This near-parity on English
together with near-zero EER on Spanish indicates that joint training achieves
bilingual generalisation, subject to the caveat in
Section~\ref{sec:datasets} that the merged corpus is not language-balanced.

\begin{table}[htbp]
\centering
\caption{Mean and std of EER (\%) for combined-trained models, Spanish eval.
         Best per back-end in \textbf{bold}; overall best
         \underline{\textbf{underlined}}.}
\label{table:mean-train-un-eval-es}
\setlength\extrarowheight{-8pt}
\footnotesize
\begin{tabular}{llcccccc}
\hline
 & & \multicolumn{2}{c}{Attention} & \multicolumn{2}{c}{Trim-pad}
   & \multicolumn{2}{c}{LSTM-sum} \\
\textbf{Front} & \textbf{Activation} &
  \textbf{Mn} & \textbf{Sd} & \textbf{Mn} & \textbf{Sd} &
  \textbf{Mn} & \textbf{Sd} \\
\hline
\textbf{CQCC}
  & AM-softmax & 1.16 & 0.43 & 1.15 & 0.38 & 0.97 & 0.24 \\
  & P2SGrad    & 0.96 & 0.36 & \textbf{0.89} & \textbf{0.25} & \textbf{0.49} & \textbf{0.18} \\
  & OC-softmax & \textbf{0.86} & \textbf{0.48} & 0.97 & 0.45 & 1.25 & 1.40 \\
  & Sigmoid    & 2.98 & 1.29 & 1.35 & 0.25 & 0.82 & 0.47 \\
\hline
\textbf{Spec.}
  & AM-softmax & 17.10 & 1.50 & 20.82 & 2.72 & 15.22 & 1.44 \\
  & P2SGrad    & \textbf{16.37} & \textbf{2.66} & 16.46 & 3.61 & 15.35 & 2.87 \\
  & OC-softmax & 17.88 & 7.00 & \textbf{17.32} & \textbf{1.62} & \textbf{14.92} & \textbf{1.83} \\
  & Sigmoid    & 16.65 & 2.40 & 18.58 & 1.35 & 16.96 & 2.74 \\
\hline
\textbf{MFCC}
  & AM-softmax & 0.31 & 0.22 & 0.42 & 0.17 & \underline{\textbf{0.09}} & \underline{\textbf{0.04}} \\
  & P2SGrad    & \textbf{0.17} & \textbf{0.09} & \textbf{0.17} & \textbf{0.09} & 0.35 & 0.33 \\
  & OC-softmax & 0.40 & 0.34 & 0.29 & 0.03 & 0.45 & 0.47 \\
  & Sigmoid    & 0.55 & 0.36 & 0.38 & 0.13 & 0.45 & 0.49 \\
\hline
\textbf{LFCC}
  & AM-softmax & 0.25 & 0.14 & \textbf{0.30} & \textbf{0.09} & 0.37 & 0.40 \\
  & P2SGrad    & \textbf{0.19} & \textbf{0.10} & 0.30 & 0.15 & 0.30 & 0.26 \\
  & OC-softmax & 0.42 & 0.36 & 0.34 & 0.13 & 0.24 & 0.12 \\
  & Sigmoid    & 0.95 & 0.87 & 0.28 & 0.11 & \textbf{0.18} & \textbf{0.14} \\
\hline
\end{tabular}
\end{table}

\begin{table}[htbp]
\centering
\caption{Mean and std of EER (\%) for combined-trained models, English eval.
         Best per back-end in \textbf{bold}; overall best
         \underline{\textbf{underlined}}.}
\label{table:mean-train-un-eval-en}
\setlength\extrarowheight{-8pt}
\footnotesize
\begin{tabular}{llcccccc}
\hline
 & & \multicolumn{2}{c}{Attention} & \multicolumn{2}{c}{Trim-pad}
   & \multicolumn{2}{c}{LSTM-sum} \\
\textbf{Front} & \textbf{Activation} &
  \textbf{Mn} & \textbf{Sd} & \textbf{Mn} & \textbf{Sd} &
  \textbf{Mn} & \textbf{Sd} \\
\hline
\textbf{CQCC}
  & AM-softmax & 6.93 & 2.54 & 6.43 & 0.40 & 6.65 & 0.71 \\
  & P2SGrad    & 5.77 & 0.33 & \textbf{5.77} & \textbf{0.11} & \textbf{5.71} & \textbf{0.31} \\
  & OC-softmax & 6.41 & 0.36 & 5.83 & 0.09 & 5.77 & 0.27 \\
  & Sigmoid    & 9.10 & 0.75 & 5.87 & 0.22 & \textbf{5.51} & \textbf{0.21} \\
\hline
\textbf{Spec.}
  & AM-softmax & 11.19 & 1.35 & 9.95 & 0.51 & 10.14 & 0.54 \\
  & P2SGrad    & 10.45 & 1.57 & \textbf{9.16} & \textbf{0.42} & \textbf{9.06} & \textbf{0.49} \\
  & OC-softmax & 10.95 & 1.23 & 10.13 & 0.50 & 9.98 & 1.23 \\
  & Sigmoid    & \textbf{10.26} & \textbf{0.74} & 10.29 & 0.72 & 10.51 & 0.99 \\
\hline
\textbf{MFCC}
  & AM-softmax & 4.96 & 1.27 & 4.43 & 0.67 & 4.62 & 0.71 \\
  & P2SGrad    & 4.07 & 0.47 & \underline{\textbf{3.64}} & \underline{\textbf{0.42}} & 4.58 & 1.02 \\
  & OC-softmax & \textbf{3.86} & \textbf{0.66} & 3.75 & 0.52 & 4.97 & 1.18 \\
  & Sigmoid    & 4.64 & 0.77 & 3.66 & 0.37 & \textbf{4.26} & \textbf{0.84} \\
\hline
\textbf{LFCC}
  & AM-softmax & 4.67 & 0.73 & 4.54 & 0.64 & 4.43 & 0.79 \\
  & P2SGrad    & \textbf{4.37} & \textbf{0.67} & 3.98 & 0.17 & \textbf{4.28} & \textbf{0.42} \\
  & OC-softmax & 5.00 & 1.31 & \textbf{3.78} & \textbf{0.44} & 4.50 & 1.22 \\
  & Sigmoid    & 5.11 & 0.74 & 3.79 & 0.49 & 4.24 & 0.51 \\
\hline
\end{tabular}
\end{table}

\subsection{Accent Analysis}
\label{sec:accent-analysis}

Using the bilingually trained models, we break down Spanish evaluation EER
by Latin-American accent.  Results below are for the MFCC front-end; LFCC
yields qualitatively similar patterns.

Table~\ref{table:accent-gender} reports the same results resolved by country
and by speaker gender, and three observations follow.  First, differences
between accents are small: Chile ($0.032\%$) and Peru ($0.053\%$) are the
best-modelled and Colombia ($0.287\%$) the worst, but all five medians are
zero and the spread is below half a percentage point, so accent alone is a
weak predictor of residual error.  Second, the back-end ranking does not vary
by accent: LCNN-LSTM-sum attains the lowest mean EER for \emph{all five}
countries, a negative result that argues against accent-conditioned model
selection rather than for it.  Third, speaker gender carries the dominant
residual signal and its direction reverses between countries: Colombia and
Venezuela are markedly worse on male speakers ($0.544\%$ against $0.029\%$,
and $0.271\%$ against $0.065\%$), whereas Argentina, Chile and Peru are worse
on female speakers.  The Colombian gap accounts for essentially all of that
country's disadvantage, and because the direction is inconsistent it cannot
reflect a single global bias of the front-end, pointing instead to an
interaction between speaker population and the attack generators used for
each subset.

\begin{table}[htbp]
\centering
\caption{Mean EER (\%) by country and speaker gender for MFCC front-end
         models trained on the combined corpus, aggregated over three
         back-ends, four losses and six seeds.}
\label{table:accent-gender}
\setlength\extrarowheight{-2pt}
\begin{tabular}{lccc}
\hline
\textbf{Country} & \textbf{Overall} & \textbf{Female} & \textbf{Male} \\
\hline
Chile      & 0.032 & 0.059 & 0.005 \\
Peru       & 0.053 & 0.085 & 0.021 \\
Argentina  & 0.140 & 0.177 & 0.103 \\
Venezuela  & 0.168 & 0.065 & 0.271 \\
Colombia   & 0.287 & 0.029 & 0.544 \\
\hline
\end{tabular}
\end{table}

\subsection{Filter-Bank Scale Effects}
\label{sec:filter-bank}

Filter-bank scale (ascendant, constant, or descendant frequency weighting)
affects primarily the cepstral front-ends.  To separate a genuine effect from
a single fortunate configuration, we evaluate all three scales across the
full grid of three back-ends and four losses, six seeds each, giving 72 runs
per scale within each filter-bank family.  Results are given in
Table~\ref{table:scale-aggregate}.

\begin{table}[htbp]
\centering
\caption{Filter-bank scale compared across the full configuration grid
         (English eval; 72 runs per scale per family).  ``Wins'' counts the
         configurations out of twelve for which that scale attains the lowest
         mean EER.}
\label{table:scale-aggregate}
\setlength\extrarowheight{-2pt}
\begin{tabular}{llccc}
\hline
\textbf{Family} & \textbf{Scale} & \textbf{Median} & \textbf{Mean} & \textbf{Wins} \\
\hline
\multirow{3}{*}{Linear (LFCC)}
   & Ascendant  & \textbf{3.069} & 3.292 & 3 \\
   & Constant   & 3.162 & 3.294 & \textbf{6} \\
   & Descendant & 3.190 & 3.383 & 3 \\
\hline
\multirow{3}{*}{Mel (MFCC)}
   & Ascendant  & 3.793 & 4.368 & 3 \\
   & Constant   & 3.855 & 4.385 & 1 \\
   & Descendant & \textbf{3.475} & \textbf{3.925} & \textbf{8} \\
\hline
\end{tabular}
\end{table}

The benefit of descendant scaling is specific to the Mel family.
For Mel-scale banks it attains the lowest median and mean and wins eight of
twelve configurations; for linear-scale banks the same weighting is the
\emph{worst} of the three by median and wins only three, with the constant
bank winning six.  Reporting either family alone would invert the conclusion,
and the assumption that low-frequency emphasis is universally preferable does
not survive the full grid.  A single configuration cannot settle this: for
LFCC with LCNN-LSTM-sum/P2S-Grad the three scales differ by less than
$0.16$~percentage points of median EER, well inside the seed-to-seed spread.

The physiology supports the restricted claim rather than the general one.  F1
does lie below roughly $1$~kHz for most vowels, but F2 spans approximately
$600$--$2500$~Hz, so an emphasis profile motivated by ``formants are low'' is
only partially justified.  The Mel scale already compresses high frequencies
and descendant weighting compounds that compression, whereas a linear bank has
no such compression and suppressing its high-frequency resolution discards
information the classifier was using.

The linear family is nonetheless the more repeatable choice, with roughly half
the across-configuration standard deviation of the Mel family under every
scale: prefer a linear bank for run-to-run stability, and a Mel bank with
descendant weighting for lowest expected EER.

%% ----------------------------------------------------------------
%\section{Conclusions}

%This paper presented a comprehensive, multi-condition evaluation of speech
%anti-spoofing countermeasures, examining the impact of language, accent, and
%filter-bank design on system performance.  The main conclusions are:

%\begin{enumerate}
%  \item \textbf{Cepstral front-ends dominate, but no single configuration
%        does.}  LFCC and MFCC consistently outperform CQCC and Spectrogram.
%        Within that group, however, the leading configurations are not
%        separable: on English none of the seven runners-up differs
%        significantly from the best configuration, and on bilingually trained
%        Spanish evaluation the four leading configurations have overlapping
 %       intervals.  Reporting a single winner would overstate what six seeds
 %       can support.

%  \item \textbf{The cross-language collapse is an attack-family effect, not a
%        uniform one.}  Models trained exclusively on English fall to chance on
%        Spanish, confirming Tamayo et~al.~\cite{tamayo2022voice}.  Decomposing
%        by attack shows the failure is concentrated in voice conversion, where
%        EER reaches $43$--$66\%$, while text-to-speech attacks remain
%        detectable at $1.79\%$ EER by the best English-trained model.  Vocoder
%        artefacts transfer across languages; conversion artefacts, entangled
%        with speaker and accent characteristics, do not.

  %\item \textbf{Bilingual training achieves near-parity on both languages.}
   %     Merging English and Spanish training data yields near-zero EER on
    %    Spanish while sacrificing fewer than one percentage point on English.
     %   We note that the merged corpus is not language-balanced
%        ($57.7\%$ Spanish), so the magnitude of both effects is confounded
 %       with corpus composition.

 % \item \textbf{Residual error tracks speaker gender, not regional accent.}
 %       Once bilingual training is in place, differences between the five
 %       Latin-American varieties are below half a percentage point and the
 %       back-end ranking is identical for all of them.  The larger and less
 %       expected effect is a gender gap whose direction reverses by country.
 %       This argues for demographic balancing rather than accent-adaptive
 %       architecture selection.

  %\item \textbf{Descendant filter-bank scaling helps Mel banks only.}  Across
   %     the full configuration grid, descendant weighting attains the lowest
    %    median EER for Mel-scale front-ends and wins eight of twelve
    %    configurations, but is the worst of the three scales for linear-scale
    %    front-ends.  Low-frequency emphasis is therefore not universally
    %    preferable. The partial justification often given for it, namely
    %    that the principal formants are low, holds for F1 but not for F2.
%\end{enumerate}

%% ----------------------------------------------------------------
\section{Conclusions}
This paper evaluated speech anti-spoofing countermeasures across languages,
accents and filter-bank designs, and five findings emerge from the results.

Cepstral front-ends outperform the alternatives, although the margin among
the best of them is too small to resolve. LFCC and MFCC place consistently
ahead of CQCC and Spectrogram, yet within that leading group the individual
configurations cannot be separated. On English, none of the seven runners-up
differs significantly from the best configuration, and under bilingual
training the four leading configurations on Spanish evaluation have
overlapping confidence intervals. Six seeds do not support naming a single
winner.

The collapse observed when the evaluation language changes is concentrated in
one attack family rather than spread evenly across all of them. Models trained
exclusively on English fall to chance level on Spanish, which confirms the
observation of Tamayo et~al.~\cite{tamayo2022voice}. Decomposing the error by
attack shows that voice conversion accounts for almost all of the degradation,
with EER between $43\%$ and $66\%$, while text-to-speech attacks remain
detectable and reach $1.79\%$ EER for the best English-trained model. Vocoder
artefacts appear to survive the change of language, whereas conversion
artefacts, entangled as they are with speaker and accent characteristics, do
not survive it.

Bilingual training reaches comparable performance in both languages. Merging
the English and Spanish training partitions drives EER on Spanish close to
zero and costs less than one percentage point on English. The merged corpus is
not balanced by language, since $57.7\%$ of its material is Spanish, so the
magnitude of both effects remains confounded with corpus composition.

The residual error that survives bilingual training follows speaker gender and
shows almost no dependence on regional accent. Differences between the five
Latin-American varieties stay below half a percentage point, and the back-end
ranking is identical for all of them. The larger and less expected effect is a
gender gap whose direction reverses from one country to another. This points
towards demographic balancing as a more productive line of work than
accent-adaptive architecture selection.

Descendant filter-bank scaling benefits Mel banks and harms linear ones.
Across the full configuration grid, descendant weighting attains the lowest
median EER for Mel-scale front-ends and wins eight of the twelve
configurations, while for linear-scale front-ends it is the worst of the three
scalings. Emphasising low frequencies is therefore not universally
advantageous. The justification usually offered for it, that the principal
formants lie in the lower part of the spectrum, holds for F1 and fails for F2.

%% ----------------------------------------------------------------
\section{Future Work}
Two limitations of this evaluation shape our ongoing work.

The first concerns attack coverage. Section~\ref{sec:cross-language} shows
that conclusions drawn across languages depend strongly on which attack
families a corpus contains, and HABLA covers only four. We are therefore
constructing \textit{MARSA}, a multi-accent Spanish corpus that adds
full-spoof attacks from six contemporary text-to-speech systems, a
partial-spoof tier in which individual words are spliced into otherwise bona
fide utterances, and an acoustic augmentation layer covering reverberation,
noise and codec degradation. Evaluating the partial-spoof tier requires
frame-level scoring, which the utterance-level EER used here cannot express.

The second concerns the training and evaluation protocol. Retraining on
language-balanced subsets would separate the effect of bilingual data from
that of corpus composition (Section~\ref{sec:datasets}), and a protocol that
also holds out the target identities used by the conversion attacks would
remove the residual optimism identified there. The gender asymmetry reported
in Section~\ref{sec:accent-analysis} further suggests that demographic
conditioning is a more promising adaptation target than regional conditioning.

Beyond these two lines, we intend to evaluate transformer-based back-ends and
self-supervised front-end representations such as
wav2vec~2.0~\cite{baevski2020wav2vec} and HuBERT~\cite{hsu2021hubert}, and to
extend the evaluation to deepfake detection with the ASVspoof~2021
dataset~\cite{delgado2021asvspoof}.




%% ----------------------------------------------------------------
%\subsubsection{Acknowledgements.}
%This work was carried out at the FLAG Lab, Universidad de los Andes.

%\subsubsection{Disclosure of Interests.}
%The authors have no competing interests to declare that are relevant to the
%content of this article.

%% ----------------------------------------------------------------
\bibliographystyle{splncs04}
\bibliography{bib/local,bib/writing}

\end{document}
